Computing the cardinality of a set intersection, $|A \cap B|$, over a shared universe of possible
elements is a core primitive in similarity search and near-duplicate detection
\cite{bayardo2007allpairs,xiao2008ppjoin}, in chemical fingerprint screening
\cite{swamidass2007bounds,haque2011tanimoto,dalke2019chemfp}, in graph co-occurrence analysis, and
in the boolean query evaluation underneath inverted and bitmap indexes. It is a single-operation
primitive: whatever the surrounding workload, what a system pays is the cost of one operation
between two sets, and that cost is fixed by how each of the two is represented before the operation
is ever issued.

The default way to compute $|X_i \cap X_j|$ represents each set as a dense bitmap over a shared
universe of $m$ possible elements and evaluates \texttt{popcount(A AND B)} with the widest
available vector instruction. This costs $\Theta(m)$ irrespective of the cardinality of either set
or of their intersection --- every one of the $m$ bits is read and combined regardless of how many
are set --- and the underlying AND-then-popcount kernel is itself a solved problem, engineered
close to the throughput limit of the hardware \cite{mula2018popcount}. For a single pair, or for
sets that occupy a substantial fraction of the universe, this is exactly the right way to compute
the answer.

Real degree and frequency distributions are rarely uniform: a small number of categories carry
most of the mass while a long tail is rare. The canonical form of this skew --- the $i$-th most
common category scaling as $\theta/i$ --- was derived for population-genetic allele-frequency
spectra \cite{fu1995} and recurs, in shape if not in constant, in item popularity, node degree, and
category-membership counts across many other kinds of data. Under such skew, row density in the
shared-universe representation spans orders of magnitude within a single corpus, so most rows are
sparse and most pairs have at least one near-empty side --- yet the $\Theta(m)$ bitmap kernel pays
the same fixed cost regardless, spending most of that cost combining zeros against zeros.

That waste is not an anecdote about one kind of data; it is a property of the fixed cost itself, and
its size can be settled before any kernel is written. Every alternative encoding of the same set ---
a sorted array of its elements, a list of its maximal runs, a fill-compressed word stream, or any of
these applied to its complement --- is exact and interconvertible with the bitmap, and each is sized
by the set rather than by the universe. The cheapest of them is smaller than the bitmap by a factor
that grows without bound as either operand moves away from density one half, so a flat cost provably
leaves headroom at every density but the middle of the range, for a single operation as much as for
a workload of them. Assuming nothing about the data beyond its density gives a floor on that
headroom rather than an estimate of it: independently placed elements are the worst case for every
encoding whose cost depends on how a set is arranged, so real structure moves those encodings
further below the bitmap and never above it. What this paper optimizes is therefore not how fast a
kernel runs but how much of the computation is never entered. The same analysis says where that is
not possible. Around density one half no encoding is smaller than the bitmap and nothing can be
skipped; there the fixed-cost kernel is the right one, and the return comes from executing it as
fast as the hardware allows --- which is what two decades of vectorized \texttt{AND}-and-popcount
engineering already deliver \cite{mula2018popcount}. Avoid the work where it can be avoided,
accelerate it where it cannot: this paper contributes the first axis, and a cheap way of knowing
which of the two applies.

Roaring is the state of the art for exactly this problem, and it already dispatches, per
fixed-width chunk, among an array, a bitset, and a run container
\cite{chambi2016roaring,lemire2018roaring} --- the same coarse-grained skip toward a cheaper
structure that inverted-index systems such as BitFunnel apply to disjointness testing
\cite{goodwin2017bitfunnel}. It is excellent, it is deployed at scale, and this paper does not
attempt to out-engineer it at its own data structure. Its adaptivity is nonetheless fixed in
granularity and static in time: which container a chunk receives is decided by a compile-time
cardinality threshold applied to that chunk alone, chosen once at ingest for a workload built
around a single pairwise operation (Methods). That same fixed commitment prices a run container
against a bitset by iterating every run and charging for its length --- a cost a
length-independent index would remove.

All of that leaves the caller's question untouched, and a caller willing to narrow it can be given
more. If only pairs above a similarity threshold are wanted, or only operands inside a band of set
sizes, then the two cardinalities alone cap the similarity a pair can reach, so most pairs are
discarded before either operand is read \cite{swamidass2007bounds,bayardo2007allpairs} --- knowing
what is being looked for is itself a way of not doing work, and it is the cheapest one available.

Two questions follow directly. Given that popcount is a solved problem and Roaring already
dispatches on container type per chunk, is there anything left for representation choice to buy?
And if a wider representation matrix does buy something, can the choice itself be made cheaply
enough --- for one operation between two large sets, and for a batch of many small ones --- that
the saving survives the deciding?

Here we treat the representation pairing itself as the object of choice. We define a pairing
matrix of \nm{count of representation-pairing cells evaluated} cells over dense-bitmap,
sorted-array, run, and fill-compressed representations and their complements, and we measure, per
cell, which pairing is fastest across the full density range, on \nm{count of microarchitectures
tested} microarchitectures, and on \nm{count of real corpora tested} real corpora drawn from graph,
text, and tabular sources. A throughput-optimized kernel wins none of \nm{count of asymmetric
measurement points tested} asymmetric measurement points tested, while selecting among
representation pairings wins on every real corpus tested, at a selection cost that stays within a
small, budgeted fraction of runtime. We also report where the approach does not pay: a
mid-density band in which the plain bitmap kernel is, and remains, the right choice.
